\documentclass[11pt,a4paper]{article}
\usepackage[margin=2.5cm]{geometry}
\usepackage{amsmath,amssymb,booktabs,xcolor,enumitem}
\usepackage[normalem]{ulem}
\usepackage[hidelinks]{hyperref}
\setlength{\parskip}{0.6em}\setlength{\parindent}{0pt}
\definecolor{newc}{HTML}{1f6f43}
\newcommand{\new}[1]{\textcolor{newc}{#1}}
\newcommand{\old}[1]{\textcolor{gray}{\sout{#1}}}
\newcommand{\oldmath}[1]{\textcolor{gray}{#1}}
\title{Proposed change [A96], the keep-up check counts the work served in the window}
\date{16 November 2026, version 1}
\begin{document}
\maketitle

\section*{How to read this}
Part~1 says what is wrong with the check as it stands. Part~2 shows the ``Keeping up'' paragraph of
the Problem Formulation as it would read. \new{Green} is new text, \old{struck} text goes, a removed
equation is grey and marked ``removed''. Part~3 has the questions. Nothing is changed in the thesis yet.

\section*{Part 1. What is wrong}

The check compares the tasks finished, $|\mathcal{C}_T|$, with the tasks finished plus the backlog
left at $T$. Since [A66] only tasks released in the measured window $[T_{\mathrm{warm}}, T]$ count in
$\mathcal{C}_T$. A task released in the warm-up and finished in the window counts nowhere.

For a fleet that keeps up this makes no difference. For a fleet that falls behind it does. Such a
fleet enters the window with a queue from the warm-up. Its agents keep working, but on those older
tasks, and the tasks released in the window wait at the back of the queue. So $|\mathcal{C}_T|$ is
close to zero even while the agents finish a task every few dozen steps.

The fleet 16 tests show it. No learned arm keeps up, and the completed share is near zero for most of
them. The centralised arm under fixed storage scores 0.13, yet on average it never goes more than 68
steps without finishing a task. The number says that the new tasks were not reached, not how much work
the fleet does. Every learned run then reads ``0 of 20'', and the arms cannot be told apart.

\textbf{The change.} Count every task finished in the window, whenever it was released, and compare
it with the tasks released in the window. The ratio is the share of the arriving work the fleet
serves. At 1 the queue holds its size. At 0.6 the fleet serves 60\,\% of what arrives, and the queue
grows by the other 40\,\%. The check then asks for a share of at least $f_{\mathrm{up}}$, for picks
and put-aways apart as before ([A59]). A run that fails it is reported by this share, which now
says how far it falls short.

\textbf{A rough estimate} from the backlog slope of the stored runs gives shares of about 0.5 to 0.7
for centralised, section-based and decentralised, and about 0.4 with messages, against about 1 for
Token Passing. The exact shares need the tests played again with the new code.

\textbf{Code.} Branch \texttt{feat/served-share}, uncommitted. \texttt{RunMetrics} gains the tasks
released and finished in the window and the share, overall and for picks and put-aways.
\texttt{keeps\_up} is not changed yet, it follows the thesis once this is approved.

\section*{Part 2. The text}

\subsection*{Problem Formulation, ``Keeping up''}

A fleet keeps up when it finishes tasks about as fast as they are released. If it falls behind,
unfinished tasks pile up, and an average over the finished ones no longer describes the run. This is
the bounded service time that \texttt{ma2017lifelong\_mapf} require of a solved instance. The queue is
the backlog $B_t$, the tasks released but not yet finished and the orders still waiting for stock.

Throughput is the number of tasks finished per timestep. It is the measure lifelong MAPF uses for
warehouses \texttt{li2021lifelong} and the pick rate that \texttt{krnjaic2022warehouse\_marl} optimise.
The check compares it with the rate at which tasks were actually released in the run, not the expected
rate, because orders arrive at random and one seed releases more than another.

\new{The check counts every task the fleet finishes in the measured steps, also one released before
them. A fleet that fell behind earlier still works through its older tasks, and that work is what
shows whether the queue grows or holds. So the measure is the share of the arriving work the fleet
serves. With 900 tasks released in the measured steps and 600 tasks finished in them, the share is
two thirds and the queue grows by 300.}

Picks and put-aways are checked separately, and a run must pass both. [unchanged]

Formally, with $\mathcal{C}_T$ the tasks completed within a horizon $T$,
throughput is $\Lambda_T = |\mathcal{C}_T| / T$. \new{Write $\mathcal{R}^{\mathrm{pick}}_T$ for the
picks released in the measured window and $\mathcal{F}^{\mathrm{pick}}_T$ for the picks finished in
it, whenever they were released, and the same for put-aways. The served share of a kind is the second
over the first, and the run passes the \emph{keep-up check} if}
\begin{equation*}
  \new{\frac{|\mathcal{F}^{\mathrm{pick}}_T|}{|\mathcal{R}^{\mathrm{pick}}_T|} \;\ge\; f_{\mathrm{up}}
  \quad\text{and}\quad
  \frac{|\mathcal{F}^{\mathrm{put}}_T|}{|\mathcal{R}^{\mathrm{put}}_T|} \;\ge\; f_{\mathrm{up}} ,}
\end{equation*}
\oldmath{\[
  \frac{|\mathcal{C}^{\mathrm{pick}}_T|}{T} \;\ge\; f_{\mathrm{up}}\, \frac{|\mathcal{C}^{\mathrm{pick}}_T| + B^{\mathrm{pick}}_T}{T}
  \quad\text{and}\quad
  \frac{|\mathcal{C}^{\mathrm{put}}_T|}{T} \;\ge\; f_{\mathrm{up}}\, \frac{|\mathcal{C}^{\mathrm{put}}_T| + B^{\mathrm{put}}_T}{T}
  \qquad\text{(removed)}
\]}
\old{where each right side is the release rate of that kind scaled by a threshold
$f_{\mathrm{up}} \le 1$ that leaves room for tasks still under way at the horizon.}
\new{where the threshold $f_{\mathrm{up}} \le 1$ leaves room for tasks still under way at the
horizon. Ship tasks count with the picks.} A run that does not keep up is reported by
\old{its throughput} \new{its served share} and is not compared on service time. The check looks at
one finite horizon, so passing it does not prove that the backlog stays bounded. [rest unchanged]

\subsection*{Notation table}
\new{$\mathcal{R}^{\mathrm{pick}}_T, \mathcal{F}^{\mathrm{pick}}_T$, and the same for put-aways. Tasks
released in the measured window, and tasks finished in it whenever released.} The rows of
$B^{\mathrm{pick}}_T, B^{\mathrm{put}}_T$ stay, since the backlog is still reported.

\section*{Part 3. Questions}

\begin{enumerate}[leftmargin=1.4em]
\item \textbf{Replace the check, or only report the share next to it?} I propose to replace it. The
  old check fails every run that fell behind in the warm-up, however much it serves later. The new one
  asks the question the check is for, whether the queue grows. A run that fell behind early but holds
  its queue would pass, with a long service time, and the backlog over the run stays reported.
\item \textbf{Keep $f_{\mathrm{up}} = 0.8$?} With the new share, 0.8 lets the queue grow by a fifth of
  what arrives. A stricter value such as 0.95 is closer to ``the queue holds''. Token Passing is at
  about 1, so either keeps its rows passing.
\item \textbf{Replay the fleet 16 tests} with the new code, about a day on the PC (or the cloud), so the
  exact shares go in the tables and the graph? Training is not touched.
\item \textbf{Do the [A66] measures change too?} No. Service time and the other window measures still
  count the tasks released in the window. Only the check changes.
\end{enumerate}

\end{document}
